\section{The log canonical and double-point cases}
\label{sec:low-cases}

We finish with two cases accessible through sectional Milnor numbers.
The double-point case reduces to a Whitney regular family of plane
curves after splitting off a square.  The remaining log canonical case,
of multiplicity three, is itself Whitney regular.
For a holomorphic germ in $n$ variables with
an isolated critical point, write $\mu^{(j)}$ for the Milnor number of its
restriction to a general $j$-dimensional linear subspace, and put
$\mu^*=(\mu^{(1)},\ldots,\mu^{(n)})$.  Thus $\mu^{(n)}=\mu$ and
$\mu^{(1)}=\operatorname{mult}_0(f)-1$.

We use two equisingularity theorems.  The equimultiplicity theorem of
Fern\'andez de Bobadilla and Pe\l ka states that a family of isolated
hypersurface singularities with constant finite Milnor number has constant
multiplicity \cite[Theorem~1.1]{FernandezBobadillaPelka2024}.  Teissier's
criterion states that, for a deformation of an isolated hypersurface
singularity with a section and smooth complement of the section over the
base, constancy of $\mu^*$ implies the Whitney conditions for that
complement and the section \cite[Chapter~II, Theorem~3.9]{Teissier1973}.
The absence of critical points off the section in our representatives
verifies the latter smoothness hypothesis.

\begin{proposition}\label{prop:low-cases}
Let $F(x,t)=f_t(x)$ be a holomorphic family in three spatial variables,
with $f_t(0)=0$, an isolated critical point at $0$ in every fiber, and
constant Milnor number.  If either $\operatorname{mult}_0(f_0)=2$ or the
normal surface germ $X_0=\{f_0=0\}$ is log canonical, then $F$ is
topologically right-trivial over a sufficiently small parameter disk.
\end{proposition}

\begin{proof}
Suppose first that $\operatorname{mult}_0(f_0)=2$.  Choose a spatial
direction in which the second derivative of $f_0$ is nonzero.  The implicit
function theorem solves the corresponding first-derivative equation as
$z=a(y,t)$, where $y\in\mathbb C^2$ and $a(0,t)=0$.  Translating $z$ by
$a(y,t)$ and taking a holomorphic square root of the resulting nonvanishing
quadratic coefficient gives biholomorphic coordinates, depending
holomorphically on $t$ and preserving the origin, in which
\[
 f_t=z^2+g_t(y).
\]
The Jacobian quotient is then the Jacobian quotient of $g_t$, so $g_t$ has
an isolated critical point and $\mu(g_t)=\mu(f_t)$.  Equimultiplicity for
this plane family gives constant $\mu^*(g_t)$, whose two entries are
$\operatorname{mult}_0(g_t)-1$ and $\mu(g_t)$.  Teissier's criterion and
Corollary~\ref{cor:whitney-right} give a joint
right-trivialization of $g_t$.  Extend its maps by the identity in $z$ and
conjugate by the coordinate changes at $0$ and at $t$.  This gives the
required right-trivialization of $f_t$, with central map the identity.

Now suppose $X_0$ is log canonical.  Adjunction and inversion of
adjunction identify this condition with log canonicity of
$(\mathbb C^3,X_0)$ \cite{Kawakita2007}; the different on the normal
Cartier hypersurface in a smooth ambient space is zero.  These
single-germ algebraic statements also apply to the analytic germ by its
analytic right-equivalence to a sufficiently high polynomial jet.
Put $m=\operatorname{mult}_0(f_0)$.  On the blowup of the origin the
exceptional plane has discrepancy $2-m$, hence $m\leq3$.  Since the
origin is a critical point, $m\geq2$, and the case $m=2$ is already settled.

For $m=3$, let $\widetilde X_0$ be the strict transform and $P\simeq
\mathbb P^2$ the exceptional plane.  The crepant boundary on the blowup
is $\widetilde X_0+P$.  Adjunction to $P$ shows that
$(P,C)$ is log canonical, where $C=\widetilde X_0|_P$ is the tangent
cubic with its scheme-theoretic multiplicities.  A boundary component of
a log canonical pair has coefficient at most one.  Thus $C$ is reduced.

Equimultiplicity keeps the order of $f_t$ equal to three.  Choose a
projective line transverse to $C$ at three distinct smooth points.  This
condition persists for the cubic initial forms of $f_t$ when $t$ is
small.  The corresponding plane sections, and hence general plane
sections, are ordinary triple points.  Their Milnor number is four:
the two quadratic initial forms of their partial derivatives have no
common projective zero, so their local intersection multiplicity is
$2\cdot2$.  Consequently
\[
 \mu^*(f_t)=(2,4,\mu(f_0))
\]
throughout a sufficiently small disk.  Teissier's criterion gives Whitney
regularity, and Corollary~\ref{cor:whitney-right} supplies the
right-trivialization.
\end{proof}


\subsection{Proof of Theorem~\ref{thm:main}}
If the central multiplicity is two, or the central surface is log canonical,
Proposition~\ref{prop:low-cases} proves the theorem.  Otherwise the central
surface is non-log-canonical.  Theorem~\ref{thm:beta-constant} proves
constancy of its logarithmic invariant.  Proposition~\ref{prop:resolved-family}
then supplies a smooth simultaneous resolution with relative normal-crossing
exceptional divisor after a finite disk base change.  The fixed-ball
hypotheses of Lemma~\ref{lem:ambient-tube} follow from
Lemmas~\ref{lem:local-family-properties} and~\ref{lem:fixed-ball-euler}.
Proposition~\ref{prop:resolution-right} gives the joint ambient right
trivialization after that base change, and Lemma~\ref{lem:remove-base-change}
descends it to the original disk.  The formulas preserve the parameter,
the origin section, and the defining function, with the identity at the
central parameter and jointly continuous inverses.  These are precisely
all the assertions of Theorem~\ref{thm:main}.\qed
